% 別行立ての数式が**文の目的語**になる形（「… is given by (1), where …」）。
%
% 論文でいちばん多い形。本文はこう取り出される:
%
%   `The estimator that we analyse in this section is and the weights are chosen
%    so that the variance is minimised over the simplex.`
%
% 式の中身は本文として取り出せない（取り出してもいけない）。取り出せなくてよいが、
% **落ちたことが読む側に伝わらなければならない**。字面は正しくても文の中身が 1 つ
% 抜けているので、正しく取れた文と同じ顔で並べてはいけない。
%
% そういう文は `\sq{...}` で書く（`\s{...}` と同じく本文 1 文だが、
% 「継ぎ目の印が付いてほしい」という印を正解に持たせる）。
%
% 直す前の実測: 印 0 件。継ぎ目の規則 2 は「またいだ先に段落の字下げがあるとき」
% しか印を付けないので、式の後に同じ段落が続くこの形には一度も当たらない。
%
% 段落の途中に**表**が入るだけの段落（最後の 1 文）は `\s{...}` のまま。
% 表は段落を切らずに流し込まれる浮きもので、文としては欠けていない。
% ここに印が付いたら誤爆。
\documentclass[10pt,twocolumn]{article}

\usepackage{mathptmx}
\usepackage[margin=0.75in]{geometry}
\usepackage{amsmath,amssymb,amsthm}
\newtheorem{lemma}{Lemma}

\newcommand{\s}[1]{#1}
% `\sq{...}` は**正解にだけ**書く印で、組版はしない（本文はすぐ下に普通に書く）。
% 式をまたいだ 1 文なので `\s{...}` のように地の文を囲む形には書けない。
% 中身は、下の段落の**式より前の行と式より後の行をそのまま空白でつないだもの**。
\newcommand{\sq}[1]{}

\title{Block Sparse Attention}
\author{Anonymous}
\date{}

\begin{document}
\maketitle

\section{Analysis}

\s{We evaluate three configurations on the standard suite and report the average accuracy across five random seeds.}

% 式が文の目的語になる 1 文目。正解は「式より前の行 ＋ 空白 ＋ 式より後の行」。
\sq{The estimator that we analyse in this section is and the weights are chosen so that the variance is minimised over the simplex.}
The estimator that we analyse in this section is
\begin{equation}
\hat{\mu} = \frac{1}{n}\sum_{i=1}^{n} w_i z_i ,
\end{equation}
and the weights are chosen so that the variance is minimised over the simplex.

\s{We now state the main guarantee of the paper.}

% 式が文の目的語になる 2 文目。
\sq{For every step size that satisfies the stated bound the iterates satisfy and the constant does not depend on the dimension of the problem.}
For every step size that satisfies the stated bound the iterates satisfy
\begin{equation}
\|x_t - x^\star\|^2 \le (1-\eta\mu)^t \|x_0 - x^\star\|^2 ,
\end{equation}
and the constant does not depend on the dimension of the problem.

\s{Under these assumptions the analysis of the previous section applies verbatim.}
\s{The kernel is implemented once and reused for both training and inference in all of the experiments that we report.}

% 段落の途中に表が入るだけの段落。表は段落を切らずに流し込まれる浮きもので、
% 文としては欠けていない。ここに印が付いたら誤爆。
\s{Training with the sparse kernel from scratch converges in the same number of steps as the dense baseline, which we verified on three seeds.}

\begin{table}[h]\centering
\begin{tabular}{lrr}
Model & Acc & F1 \\
Dense & 81.2 & 79.4 \\
Sparse & 83.7 & 81.9 \\
\end{tabular}
\caption{Accuracy and F1 on the held-out split.}
\end{table}

\end{document}
